% =============================================================================
% APENDICE B: DETALLES DE INGENIERIA DE CARACTERISTICAS
% (fusion del antiguo apendice-lista con la descripcion detallada, julio 2026)
% =============================================================================

\section{Detalles de Ingenier\'ia de Caracter\'isticas}
\label{sec:appendix_features}

Las 536 caracter\'isticas del sistema se derivan de las 92 variables base mediante ocho familias de transformaci\'on; a continuaci\'on se describe cada una, su motivaci\'on y las caracter\'isticas que produce.

\paragraph{Indicadores t\'ecnicos cl\'asicos (19 caracter\'isticas).}
Los precios en bruto (apertura, m\'aximo, m\'inimo, cierre y volumen) contienen informaci\'on que no es directamente accesible para un modelo predictivo en su forma original. Un precio de cierre de \$450 no indica por s\'i solo si el mercado est\'a sobrecomprado, si la tendencia se est\'a agotando o si el volumen confirma el movimiento. Los indicadores t\'ecnicos resuelven esta limitaci\'on condensando la historia reciente de precio y volumen en se\~nales interpretables que representan el estado actual del mercado. Para un modelo de \textit{machine learning}, estos indicadores no funcionan como reglas de decisi\'on (``comprar si RSI $<$ 30''), sino como representaciones compactas del r\'egimen de precio sobre las cuales el modelo puede aprender qu\'e combinaciones de condiciones preceden retornos positivos o negativos.

A partir de los datos OHLCV del SPY y utilizando la librer\'ia \hyperlink{glos:talib}{TA-Lib}, se calculan el \'Indice de Fuerza Relativa (RSI) en ventanas de 7, 14 y 21 d\'ias, que mide la magnitud relativa de ganancias frente a p\'erdidas recientes; la Convergencia/Divergencia de Medias M\'oviles (MACD) con sus componentes de l\'inea, se\~nal e histograma; el \'Indice Direccional Medio (ADX) junto con los indicadores direccionales +DI y $-$DI, que cuantifican la fuerza y direcci\'on de la tendencia; el Oscilador Estoc\'astico (\%K, \%D) y el Williams \%R, que identifican condiciones de sobrecompra y sobreventa; el \'Indice de Canal de \textit{Commodities} (CCI); el Rango Verdadero Medio (ATR), que mide la volatilidad intrad\'ia; el \textit{On-Balance Volume} (OBV), que acumula volumen ponderado por la direcci\'on del precio; y la posici\'on del precio dentro de las Bandas de Bollinger. Estos 19 indicadores constituyen el vocabulario base del an\'alisis t\'ecnico sobre el cual se construyen las familias m\'as especializadas que se describen a continuaci\'on.

\paragraph{\textit{Momentum}, retornos y volumen (32 caracter\'isticas).}
Mientras que los indicadores t\'ecnicos cl\'asicos describen el \textit{estado} actual del mercado, las caracter\'isticas de \textit{momentum} capturan su \textit{inercia}, es decir, la velocidad y direcci\'on del cambio de precio a distintas escalas temporales. Esta distinci\'on es relevante porque la evidencia emp\'irica muestra que el \hyperlink{glos:momentum}{\textit{momentum}} opera de forma diferente seg\'un el horizonte, dado que a muy corto plazo (1--5 d\'ias) los retornos tienden a revertir, a mediano plazo (1--12 meses) tienden a persistir, y a largo plazo (3--5 a\~nos) vuelven a revertir. Proveer m\'ultiples ventanas permite a los modelos aprender cu\'al de estas din\'amicas domina en cada r\'egimen de mercado.

Los retornos hist\'oricos se calculan mediante la funci\'on \texttt{pct\_change}$(n)$, que computa el cambio porcentual entre el precio actual y el precio de hace $n$ d\'ias h\'abiles, dado por $r_n[t] = (P[t] - P[t-n]) / P[t-n]$. Se generan retornos para ventanas de 1, 5, 21, 63, 126 y 252 d\'ias, junto con sus equivalentes expresados como tasa de cambio porcentual (\textit{Rate of Change}). La distancia del precio a medias m\'oviles de 10, 21, 50, 100 y 200 d\'ias captura la posici\'on relativa del precio respecto a su tendencia de largo plazo. Los cruces de medias m\'oviles (\textit{crossovers}), es decir, se\~nales binarias que indican cu\'ando una media m\'ovil de corto plazo supera o cae por debajo de una de largo plazo (pares 5/10, 10/21, 20/50, 50/200), codifican cambios de tendencia. La volatilidad realizada se estima como la desviaci\'on est\'andar m\'ovil de retornos diarios, anualizada por $\sqrt{252}$, en ventanas de 5 a 126 d\'ias, junto con ratios de volatilidad entre ventanas (5/21 y 21/63) que detectan compresiones o expansiones de volatilidad. Completan esta familia dos caracter\'isticas de volumen, a saber, el volumen relativo respecto a su media m\'ovil de 20 y 63 d\'ias, que mide si la actividad de mercado reciente es at\'ipicamente alta o baja.

\paragraph{An\'alisis t\'ecnico avanzado (60 caracter\'isticas).}
Los indicadores cl\'asicos descritos anteriormente miden una sola dimensi\'on del mercado a la vez, dado que el RSI captura \textit{momentum}, el ATR captura volatilidad y el OBV captura volumen. Sin embargo, las decisiones de inversi\'on rara vez dependen de una \'unica dimensi\'on, y lo relevante es la \textit{confluencia} de m\'ultiples condiciones simult\'aneas y el \textit{contexto} en que se producen (posici\'on del precio dentro de un rango hist\'orico, zona de soporte o resistencia, alineaci\'on de m\'ultiples se\~nales). Las seis familias que siguen capturan precisamente estas relaciones multidimensionales.

La \hyperlink{glos:ichimoku}{\textbf{Nube de Ichimoku}} genera 13 caracter\'isticas a partir de cinco l\'ineas calculadas sobre ventanas de 9, 26 y 52 d\'ias. La ``nube'' (espacio entre las l\'ineas Senkou A y B) define zonas de soporte y resistencia din\'amicas, y a partir de ella se codifica la posici\'on del precio respecto a la nube (por encima, dentro o por debajo), el espesor de la nube como proporci\'on del precio, las distancias a las l\'ineas de referencia (Tenkan y Kijun) y los cruces entre ellas, que se\~nalan cambios de tendencia. Una variable compuesta indica la confluencia alcista plena, es decir, precio sobre la nube y l\'inea r\'apida sobre la lenta.

Los \hyperlink{glos:fibonacci}{\textbf{retrocesos de Fibonacci}} producen 18 caracter\'isticas para ventanas de 63 y 126 d\'ias. En cada d\'ia se identifica el rango (m\'aximo y m\'inimo del per\'iodo) y se calcula la posici\'on relativa del precio dentro de ese rango, su distancia a cada nivel de Fibonacci y variables binarias que indican si el precio se encuentra en una zona de proximidad ($\pm 2\%$) a los niveles 38.2\%, 50\% y 61.8\%, donde las reversiones son estad\'isticamente m\'as frecuentes.

Los \hyperlink{glos:donchian}{\textbf{canales de Donchian}} aportan 8 caracter\'isticas para per\'iodos de 20 y 55 d\'ias, a saber, la posici\'on del precio dentro del canal, su ancho como proporci\'on del precio y se\~nales binarias de ruptura al alza o a la baja. Los \textbf{canales de Keltner} complementan con 3 caracter\'isticas que miden la posici\'on del precio dentro de bandas basadas en el ATR, e incluyen un indicador de \hyperlink{glos:squeeze}{\textit{squeeze}} que detecta cuando las Bandas de Bollinger se comprimen dentro de los canales de Keltner, condici\'on que frecuentemente anticipa expansiones de volatilidad.

El \textbf{flujo de dinero} se cuantifica mediante 8 caracter\'isticas, a saber, el \textit{Money Flow Index} (MFI) en ventanas de 14 y 21 d\'ias, un oscilador que pondera el volumen por la direcci\'on del precio \hyperlink{glos:precio_tipico}{t\'ipico} con se\~nales de sobrecompra y sobreventa, y el Chaikin \textit{Money Flow} (CMF) en ventanas de 20 y 50 d\'ias, que mide la presi\'on acumulada de compra o venta seg\'un d\'onde cierra el precio dentro del rango intrad\'ia. Por \'ultimo, 10 \textbf{se\~nales de confluencia estrat\'egica} codifican combinaciones binarias de condiciones de mercado, incluyendo estados del RSI, posici\'on del precio respecto a medias m\'oviles de largo plazo (SMA-200), cruces dorados y de muerte (SMA-50 vs.\ SMA-200), fuerza de tendencia (ADX $>$ 25) y puntuaciones compuestas que requieren la alineaci\'on simult\'anea de al menos tres condiciones alcistas o bajistas.

\paragraph{An\'alisis multitemporal Triple Screen (24 caracter\'isticas).}
Se implementa el sistema Triple Screen de \citet{elder1993}, adaptado como generador de caracter\'isticas; a diferencia de su uso como estrategia de \textit{trading}, en este estudio se adopta exclusivamente como fuente de variables predictivas. El principio central de este sistema es que una decisi\'on de inversi\'on mejora cuando se analiza el mercado en m\'ultiples escalas temporales simult\'aneamente, dado que primero se identifica la tendencia dominante en un horizonte largo, luego se buscan puntos de entrada en un horizonte intermedio, y finalmente se confirma la se\~nal con indicadores de corto plazo. Elder denomin\'o a estos tres niveles de an\'alisis \textit{screens} (filtros), de ah\'i el nombre del sistema.

En el primer nivel, se simula un horizonte semanal a partir de datos diarios mediante ventanas m\'oviles de cinco d\'ias, calculando un MACD con per\'iodos expandidos (multiplicados por 5) para identificar la direcci\'on de la tendencia de fondo. En el segundo nivel, se generan osciladores diarios que miden presiones de corto plazo \textit{dentro} de esa tendencia, a saber, el \'Indice de Fuerza de Elder en versiones de 2 y 13 per\'iodos, que combina cambio de precio y volumen normalizado por su desviaci\'on est\'andar m\'ovil de 63 d\'ias, y el \textit{Elder Ray} (\textit{Bull/Bear Power}), que mide la distancia entre los extremos intrad\'ia y la media m\'ovil exponencial (EMA) de 13 per\'iodos. En el tercer nivel, el Sistema de Impulso combina la direcci\'on de la EMA con la del histograma MACD para clasificar cada d\'ia como alcista (ambos subiendo), bajista (ambos cayendo) o neutral. La l\'ogica subyacente es que las se\~nales m\'as fiables surgen cuando los tres niveles est\'an alineados, es decir, una tendencia alcista de fondo, un retroceso temporal en los osciladores y una confirmaci\'on de impulso diario. Se incluyen adem\'as se\~nales puras de Elder, a saber, compra cuando \textit{Bear Power} es negativo pero creciente (anticipaci\'on de giro alcista) y venta cuando \textit{Bull Power} es positivo pero decreciente, y un \textit{score} compuesto que cuantifica el grado de alineaci\'on entre los tres niveles.

\paragraph{Estimadores de volatilidad OHLC y microestructura (81 caracter\'isticas).}
La volatilidad calculada a partir de retornos \textit{close-to-close}, es decir, la desviaci\'on est\'andar m\'ovil de los retornos diarios ya incluida en la familia de \textit{momentum}, presenta una limitaci\'on importante, ya que descarta la informaci\'on contenida en el movimiento intrad\'ia del precio. Un d\'ia en que el mercado sube un 3\% y luego retrocede para cerrar sin cambios registra retorno cero y, por tanto, volatilidad \textit{close-to-close} nula, a pesar de haber experimentado una dispersi\'on de precio considerable. Los datos OHLC (apertura, m\'aximo, m\'inimo y cierre) capturan esta informaci\'on perdida, y la literatura financiera ha producido una secuencia de estimadores progresivamente m\'as eficientes para explotarla.

\citet{parkinson1980} propone el primer estimador basado exclusivamente en el rango \textit{High-Low}, sugiriendo que es aproximadamente cinco veces m\'as eficiente estad\'isticamente que la volatilidad \textit{close-to-close}. \citet{garman1980} mejoran esta eficiencia incorporando la relaci\'on \textit{Close-Open}, que a\~nade informaci\'on sobre la direcci\'on del movimiento intrad\'ia. Ambos estimadores asumen, sin embargo, que la serie no presenta tendencia (\hyperlink{glos:drift}{\textit{drift}}), supuesto frecuentemente violado en mercados accionarios, donde los retornos esperados de largo plazo son positivos. \citet{rogers1991} eliminan esta limitaci\'on con un estimador que separa expl\'icitamente el \textit{drift} de la volatilidad. Finalmente, \citet{yang2000} integran los aportes anteriores en un estimador que descompone la varianza total en tres componentes, a saber, \textit{overnight} (el \textit{gap} entre cierre anterior y apertura), \textit{open-to-close} (movimiento intrad\'ia direccional) y Rogers-Satchell (dispersi\'on pura), combinados con una ponderaci\'on \'optima ($k = 0.34$) que minimiza la varianza del estimador. Cada uno de estos cuatro estimadores se calcula en ventanas de 21 y 63 d\'ias, produciendo 8 caracter\'isticas cl\'asicas.

Un conjunto posterior de trabajos no se limita a estimar la volatilidad puntual, sino que busca extraer caracter\'isticas adicionales de la din\'amica del rango. \citet{alizadeh2002} demuestran que el logaritmo del rango, $\ln(H/L)$, sigue una distribuci\'on aproximadamente gaussiana (a diferencia de la volatilidad misma, que es marcadamente asim\'etrica), lo que permite aplicar herramientas estad\'isticas est\'andar (medias, desviaciones, percentiles, momentos de orden superior) directamente sobre esta variable. \citet{chou2005} modela el rango diario como un proceso autorregresivo condicional (\hyperlink{glos:carr}{CARR}), capturando su persistencia temporal mediante rezagos, innovaciones (rango observado menos rango esperado) e indicadores de \textit{shock} cuando el rango excede dos desviaciones est\'andar de su distribuci\'on reciente. \citet{fiszeder2016} combinan retornos y rango en un marco unificado (\textit{Range-GARCH}), aportando medidas de eficiencia de precio (qu\'e fracci\'on del rango intrad\'ia se traduce en retorno direccional) y descomposiciones de la varianza en componentes \textit{overnight} e intrad\'ia. \citet{vinte2021} aplican teor\'ia de informaci\'on al an\'alisis de la dispersi\'on intrad\'ia, midiendo la entrop\'ia de la posici\'on del precio dentro del rango diario ponderada por el volumen negociado; un d\'ia donde apertura y cierre coinciden con volumen elevado (alta entrop\'ia) refleja indecisi\'on activa del mercado, mientras que uno donde el precio recorre el rango completo (baja entrop\'ia) refleja convicci\'on direccional. En conjunto, estos estimadores contempor\'aneos generan 54 caracter\'isticas.

Complementan esta dimensi\'on 19 caracter\'isticas de \textbf{microestructura} que diseccionan la anatom\'ia de cada sesi\'on individual, incluyendo el rango intrad\'ia y el \hyperlink{glos:true_range}{\textit{True Range}}, la posici\'on del cierre dentro del rango diario (0 = cerr\'o en el m\'inimo, 1 = cerr\'o en el m\'aximo), la magnitud y direcci\'on de los \hyperlink{glos:gap_overnight}{\textit{gaps overnight}} junto con un indicador de si el \textit{gap} fue ``llenado'' durante la sesi\'on, patrones de velas japonesas (proporci\'on de cuerpo, sombras superior e inferior, detecci\'on de \hyperlink{glos:doji}{doji}), volumen relativo respecto a promedios de 21 y 63 d\'ias, y el producto rango-volumen como indicador compuesto de actividad de mercado.

\paragraph{Relaciones \textit{cross-asset}, VIX y condiciones macro (66 caracter\'isticas).}
Las familias anteriores analizan el S\&P~500 en aislamiento, considerando su precio, su volumen y su volatilidad. Sin embargo, el mercado accionario no se mueve de forma independiente, sino que reacciona a flujos de capital entre clases de activos, a cambios en las expectativas de pol\'itica monetaria y al estado del ciclo econ\'omico real. La literatura ha documentado ampliamente estos v\'inculos; \citet{chen1986} demuestran que variables macroecon\'omicas como la producci\'on industrial, la inflaci\'on y la estructura de tasas de inter\'es explican una parte significativa de los retornos accionarios; \citet{fama1989} muestran que los \textit{spreads} de cr\'edito y el \textit{dividend yield} predicen retornos del mercado a horizontes de uno a cuatro a\~nos; \citet{baker2006} documentan que el sentimiento del inversor afecta la secci\'on transversal de retornos esperados; y \citet{hong2012} encuentran que la actividad en los mercados de \textit{commodities} anticipa movimientos en el mercado accionario. Sin estas variables, el modelo solo observa el precio del SPY; con ellas, dispone del contexto de otras clases de activos que podr\'ian influir en su comportamiento.

Las relaciones \hyperlink{glos:cross_asset}{\textit{cross-asset}} se cuantifican mediante correlaciones m\'oviles entre los retornos del SPY y cinco activos de referencia (bonos del Tesoro a 10 a\~nos, oro, petr\'oleo, d\'olar y VIX) en ventanas de 21 y 63 d\'ias. Estas correlaciones son din\'amicas, dado que en per\'iodos de \textit{risk-off} acciones y bonos se descorrelacionan (los inversores venden renta variable y compran deuda soberana como refugio), mientras que en per\'iodos de \textit{risk-on} la correlaci\'on se estrecha. Las caracter\'isticas del VIX incluyen cambios a 1, 5 y 21 d\'ias, indicadores de r\'egimen basados en percentiles hist\'oricos (percentil 75 y 25 sobre una ventana m\'ovil de 252 d\'ias), la estructura temporal (en particular la relaci\'on VIX3M/VIX, donde un ratio mayor a 1 indica \hyperlink{glos:contango}{\textit{contango}} y un ratio menor a 1 indica \textit{backwardation}) y la prima de riesgo de volatilidad, definida como la diferencia entre la volatilidad impl\'icita y la volatilidad realizada.

Se incorporan caracter\'isticas de la curva de rendimiento, el predictor macroecon\'omico de recesiones m\'as documentado en la literatura \citep{harvey1989}, incluyendo el \textit{spread} 10Y--2Y y 10Y--3M, indicadores de inversi\'on y el \hyperlink{glos:butterfly}{\textit{butterfly}} 2--5--10, junto con \textit{spreads} de cr\'edito (\textit{investment grade} y \textit{high yield}) cuya ampliaci\'on anticipa deterioro en las condiciones financieras. Los indicadores macroecon\'omicos se transforman en caracter\'isticas de \textit{momentum} (cambio a 21 y 63 d\'ias) y tendencia (cruce de media m\'ovil de 5 contra 21 d\'ias) para siete indicadores clave (PIB, PMI, empleo, desempleo, IPC, producci\'on industrial y confianza del consumidor). Completan esta familia cuatro indicadores compuestos de r\'egimen de mercado (basados en precio vs.\ MA-200, \textit{momentum}, nivel del VIX y proporci\'on de sectores positivos) y t\'erminos de interacci\'on entre variables (VIX $\times$ curva de rendimiento, \textit{momentum} condicionado a r\'egimen de volatilidad).

\paragraph{Variables rezagadas y estad\'isticos m\'oviles (78 caracter\'isticas).}
Esta familia aplica dos tipos de transformaciones dise\~nadas para enriquecer la representaci\'on temporal de las variables y facilitar el aprendizaje de los modelos.

La primera transformaci\'on genera \textbf{variables rezagadas} (\textit{lags}) que proporcionan al modelo memoria expl\'icita del estado pasado de las variables. Formalmente, un \textit{lag} de orden $n$ desplaza la serie temporal $n$ d\'ias hacia el pasado, dado por $X_{\text{lag}(n)}[t] = X[t-n]$. Por ejemplo, \texttt{VIX\_lag5}$[t]$ contiene el nivel del VIX de hace cinco d\'ias h\'abiles. Se generan \textit{lags} para 18 variables clave, entre ellas tasas de inter\'es a lo largo de la curva (Fed Funds, T-Bills, bonos a 2, 5, 10 y 30 a\~nos), VIX, \'indices de mercado (Russell 2000, EFA, EEM, XLF, XLU) e indicadores econ\'omicos (PIB, NFP, ventas minoristas), en tres horizontes (1, 5 y 21 d\'ias), totalizando 54 caracter\'isticas. Estos \textit{lags} permiten a los modelos detectar patrones de cambio temporal; por ejemplo, si el VIX se dispar\'o hace 5 d\'ias pero ha revertido a su nivel actual, esa trayectoria contiene informaci\'on predictiva que el valor contempor\'aneo por s\'i solo no captura.

La segunda transformaci\'on calcula \textbf{estad\'isticos m\'oviles} (\textit{rolling statistics}) mediante ventanas m\'oviles (\textit{rolling windows}) de 21 y 63 d\'ias. Para cada d\'ia $t$ y cada ventana de tama\~no $w$, se computa la media $\bar{X}_w[t] = \frac{1}{w}\sum_{i=0}^{w-1} X[t-i]$ y la desviaci\'on est\'andar $s_w[t]$ sobre las \'ultimas $w$ observaciones. La media m\'ovil suaviza las fluctuaciones diarias y revela la tendencia subyacente de la variable; por ejemplo, una media m\'ovil de 63 d\'ias del VIX indica su nivel t\'ipico del \'ultimo trimestre, permitiendo distinguir si el valor actual est\'a por encima o por debajo de lo habitual. La desviaci\'on est\'andar m\'ovil cuantifica cu\'an estable o vol\'atil ha sido la variable en el per\'iodo reciente. Se aplican estas transformaciones a un subconjunto de variables de referencia, generando 24 caracter\'isticas.

\paragraph{Transformaciones adicionales y factores (176 caracter\'isticas).}
Las variables en bruto descritas en la Secci\'on~3.1.1 capturan \textit{niveles}, como el precio del oro hoy, la tasa de desempleo este mes o el VIX esta ma\~nana. Sin embargo, los modelos predictivos aprenden con mayor facilidad a partir de \textit{cambios} que de niveles, por dos razones. En primer lugar, los cambios son m\'as estacionarios que los niveles (un retorno del 2\% tiene el mismo significado en 2005 que en 2025, mientras que un precio de \$150 no). En segundo lugar, los cambios capturan la din\'amica reciente de la variable, que es la informaci\'on predictivamente m\'as relevante para horizontes cortos.

Todas las variables base de inter\'es reciben un conjunto uniforme de dos transformaciones, a saber, cambio porcentual diario (\texttt{pct\_change}) y \textit{momentum} a 5 d\'ias (diferencia entre el valor actual y el de hace 5 d\'ias). Se a\~naden caracter\'isticas de factores de estilo, inspiradas en la literatura de \hyperlink{glos:factores_estilo}{\textit{asset pricing}}. Estas incluyen el \textit{spread} valor/crecimiento (diferencial de retornos entre XLF y XLK como \hyperlink{glos:proxy}{\textit{proxies}} de \textit{value} y \textit{growth}), tama\~no (\textit{small vs.\ large cap}, IWM vs.\ SPY) y \textit{momentum} tecnol\'ogico (rendimiento relativo de XLK frente al mercado amplio). Completan esta familia caracter\'isticas adicionales del sistema Elder y extensiones de las variables de VIX y curva de rendimiento.
